\chapter{Interrogazioni Nidificate}
\section{Cosa sono?}
Una delle ragioni che rendono SQL un linguaggio potente è la possibilità di esprimere Interrogazioni più complesse in termini di interrogazioni più semplici, tramite il meccanismo delle subqueries

\section{Come funzionano?}
La clauso WHERE di una query può contenere un'altra query.

La subquery viene usata per determinare uno o più valori da usare come valori di confronto in un predicato della query esterna

\section{Sintassi Subquery}
\begin{lstlisting}[language=sql_generic]
SELECT <lista_attributi>
FROM <tabella>
WHERE <valore_scalare> <operatore_confronto> [ANY | ALL] 
    (   SELECT <colonna> 
        FROM <tabella_sottoquery> 
        [WHERE <condizioni>] 
    );
\end{lstlisting}

\section{Predicati ANY e ALL}
Nella clausola WHERE possno comparire predicati che confrontano un attributo (o un'espressione sugli attributi) con il risultato di una query SQL

\begin{itemize}
    \item \textbf{ANY:} Il predicato è vero se almeno una riga restituita dalla query SelectSQL soddisfa il controllo
    \item \textbf{ALL:} Il predicato è vero se tutte le righe restituite dalla query SelectSQL soddisfano il confonto
\end{itemize}

\subsection{Sintassi}
\begin{lstlisting}[language=sql_generic]
ScalarValue Operator <ANY | ALL> SelectSQL
\end{lstlisting}

L'Operato può essere uno tra : =, <>, <, <=, >, >=

\subsection{Esempio}
\subsubsection{ANY}
\begin{lstlisting}[language=sql_generic]
SELECT Cognome  
FROM Impiegato 
WHERE Dipart = ANY (SELECT Nome
                    FROM Dip            
                    WHERE Citta = 'Milano');
\end{lstlisting}

\subsubsection{ALL}
\begin{lstlisting}[language=sql_generic]
SELECT Nome, Cognome 
FROM Impiegato 
WHERE Stipendio >= ALL (SELECT Stipendio 
                        FROM Impiegato);
\end{lstlisting}

\section{Quantificatori esistenziali}
EXISTS e NOT EXISTS ammettono come parametro le query nidificate.

\subsection{EXISTS}
L'operatore logico EXISTS restituisce il valore booleano true se la subquery restituisce almeno una tupla, false altrimenti

\subsubsection{Sintassi}
\begin{lstlisting}[language=sql_generic]
EXISTS SelectStar
\end{lstlisting}

\subsubsection{Esempio}
\begin{lstlisting}[language=sql_generic]
SELECT A.Citta, A.Nazione
FROM AEROPORTO AS A
WHERE EXISTS (SELECT * FROM VOLO AS V 
              WHERE V.CittaPart = A.Citta);
\end{lstlisting}

\subsection{NOT EXISTS}
L'operatore NOT EXISTS restituisce il valore booleano true se la subquery non restituisce alcuna tupla, false altrimenti

\subsubsection{Sintassi}
\begin{lstlisting}[language=sql_generic]
NOT EXISTS SelectStar
\end{lstlisting}

\subsubsection{Esempio}
\begin{lstlisting}[language=sql_generic]
SELECT O.CodOrd
FROM Ordine AS O
WHERE NOT EXISTS (
    SELECT * FROM Prodotto AS P
    WHERE NOT EXISTS (
        SELECT * FROM Dettaglio AS D
        WHERE D.CodOrd = O.CodOrd 
          AND D.CodProd = P.CodProd
    )
);
\end{lstlisting}